% Source: https://artofproblemsolving.com/community/c3h2604975_coeficcent_in_sfft
% Stable identity: aops-topic-2604975-q01
\textbf{(a)} Multiplying by $k\ne0$ (which is reversible) gives the equivalent equation $k^2ab+kma+knb=kR$. Since
\[(ka+n)(kb+m)=k^2ab+kma+knb+mn,\]
this is the same as $(ka+n)(kb+m)=kR+mn$. So the trick is: first multiply through by the coefficient of $ab$.

\textbf{(b)} Here $k=3$, $m=2$, $n=4$, so the equation becomes
\[(3a+4)(3b+2)=3\cdot44+8=140.\]
For positive $a,b$ we need $3a+4\ge7$ with $3a+4\equiv1\pmod3$, and $3b+2\ge5$ with $3b+2\equiv2\pmod3$. The divisors of $140$ are $1,2,4,5,7,10,14,20,28,35,70,140$. Those that are $\equiv1\pmod3$ and at least $7$ are $7,10,28,70$, with cofactors $20,14,5,2$. The cofactor must be $\equiv2\pmod3$ and at least $5$, which rules out only $70\cdot2$. So
\[3a+4=7,10,28\quad\text{and}\quad3b+2=20,14,5,\]
giving $\boxed{(a,b)=(1,6),\ (2,4),\ (8,1)}$. Check: $18+2+24$, $24+4+16$ and $24+16+4$ all equal $44$.
